\hypertarget{appendix-49d27bc9d9a5}{}
\section*{One Definition Rule}
\section{One Definition Rule}\label{appendix-49d27bc9d9a5-one-definition-rule}

\textbf{Curriculum depth:} Introduced · \textbf{First seen:} C007 · \textbf{Latest development:} C007

\subsection{Definition and mechanism}\label{appendix-49d27bc9d9a5-definition-and-mechanism}

ODR requires coherent definitions. A non-inline externally linked function normally needs one program definition when odr-used. Certain class, template and inline definitions can appear in multiple translation units only under the rule's conditions, including matching tokens and lookup; named-module attachment changes the permitted cases. Header guards operate within one translation unit, not across the program.

\subsection{Prerequisites}\label{appendix-49d27bc9d9a5-prerequisites}

\begin{itemize}
\tightlist
\item
  \hyperref[appendix-2e4d494c2d94]{Declarations and definitions} --- Agreement about which declarations define entities is necessary to state the one-definition constraint.
\end{itemize}

\subsection{Used by}\label{appendix-49d27bc9d9a5-used-by}

Later chunks extend this same node; dependent concepts are listed here as they are integrated.

\subsection{Recall and next depth}\label{appendix-49d27bc9d9a5-recall-and-next-depth}

Give an example, explain why the rule exists, and identify which guarantee would be missing if the prerequisite were ignored.

\textbf{Pending:} Deeper ownership, pointer/reference, generic programming and module rules from later sources. This is written coverage, not demonstrated personal mastery.

\subsection{Source and navigation}\label{appendix-49d27bc9d9a5-source-and-navigation}

\hyperref[01-foundations-i-ch002-many-files-must-describe-one-coherent-program]{C007 teaching chapter} · \hyperref[appendix-fd090d982a15-m007---parts-101-112]{Source coverage} · \href{../02-Foundations-II/main.pdf\#nameddest=appendix-4ad47bdd269a}{Dependency Map (02-Foundations-II, appendix-4ad47bdd269a)}
